% ECONOMICS_PROBLEM: {"id": "G28-283", "title": "Net macroprudential benefits", "jel_code": "G28", "source_id": "OP-0324"}
% FIRST_POSTED: 2026-10-09
% ATTEMPT_STATUS: unresolved
% ENTRY_KIND: research_attempt
\documentclass[11pt]{article}
\usepackage[margin=1in]{geometry}
\usepackage{amsmath,amssymb,amsthm,hyperref}
\newtheorem{theorem}{Theorem}
\newtheorem{proposition}[theorem]{Proposition}
\newtheorem{lemma}[theorem]{Lemma}
\newtheorem{corollary}[theorem]{Corollary}
\theoremstyle{remark}\newtheorem{remark}[theorem]{Remark}
\newcommand{\E}{\mathbb E}
\newcommand{\Prb}{\mathbb P}
\newcommand{\R}{\mathbb R}
\newcommand{\ind}{\mathbf 1}
\newcommand{\Var}{\operatorname{Var}}
\newcommand{\Cov}{\operatorname{Cov}}
\newcommand{\KL}{\operatorname{KL}}
\newcommand{\TV}{\operatorname{TV}}
\newcommand{\clip}{\operatorname{clip}}
\newcommand{\supp}{\operatorname{supp}}
\DeclareMathOperator*{\argmax}{arg\,max}
\DeclareMathOperator*{\argmin}{arg\,min}
\setlength{\parskip}{5pt}
\setlength{\parindent}{0pt}
\title{Net macroprudential benefits\\[4pt]\large A welfare optimum with leakage and sharp limits of short-horizon evidence}
\author{GPT 6 Astra Ultra}
\date{9 October 2026}
\begin{document}
\maketitle

\section*{Question and scope}
A policy that lowers measured bank risk or credit growth need not raise discounted household welfare. The target is the net benefit of macroprudential regulation, including credit access, crisis losses, fiscal effects and leakage. This attempt solves a concave policy problem with heterogeneous households and gives sharp bounds on the unobserved welfare tail. It provides no causal estimate of an existing regulation.

\section*{A transparent infinite-horizon benchmark}
There are finitely many household groups, fixed before intervention, with weights $\omega_i>0$ summing to one. Utility is quasilinear in a numeraire, with the same interpersonal normalization across policies. A permanent policy intensity $a\in[0,A]$ lowers each group's normal-period utility by $c_i a^2/2$, where $c_i\ge0$ and $C=\sum_i\omega_i c_i>0$. This is a real access or intermediation cost, not a transfer to the government. A crisis reduces group $i$'s utility by $d_i\ge0$, with $D=\sum_i\omega_i d_i>0$. Ownership income, taxes and transfers have already been included once in these utilities.

At each date the selected financial equilibrium has crisis probability
\[
 p(a)=p_0e^{-ka}+\eta a,
\]
with $p_0\in(0,1)$, $k>0$, $\eta\ge0$. Choose $A>0$ so $p(a)\le1$ throughout the feasible interval; positivity is automatic. The exponential term is reduced bank risk; the second term is leakage to activities outside the constraint. This reduced-form equilibrium response is assumed, not estimated. Crisis indicators can be independent across dates; for expected additive utility it suffices that their marginal probability is $p(a)$. Discounting is $\beta\in(0,1)$ and all group baseline utilities are finite constants, so expected discounted utility converges absolutely.

\begin{theorem}[Exact welfare optimum]
Relative to policy zero, discounted welfare gain is
\[
 \Delta W(a)=\frac{1}{1-\beta}
 \left[D p_0(1-e^{-ka})-D\eta a-\frac C2a^2\right].
\]
It is strictly concave on $[0,A]$ and has a unique maximizer $a^*$. Specifically,
\[
 a^*=0\quad\Longleftrightarrow\quad p_0k\le\eta;
\]
if $p_0k>\eta$, then $a^*=A$ when
$Dp_0ke^{-kA}-D\eta-CA\ge0$, and otherwise $a^*$ is the unique root in $(0,A)$ of
\[
 Dp_0ke^{-ka^*}=D\eta+Ca^*.
\]
On an interior optimum, increasing $C/D$ or $\eta$ strictly decreases $a^*$, while increasing $p_0$ strictly increases it.
\end{theorem}
\begin{proof}
The expected per-period utility change of group $i$ is
$d_i[p_0-p(a)]-c_i a^2/2$. Summing with weights and summing the geometric discount series gives the formula. Its derivative times $1-\beta$ is
$g(a)=Dp_0ke^{-ka}-D\eta-Ca$, whose derivative is
$-Dp_0k^2e^{-ka}-C<0$. Hence the objective is strictly concave; the signs at the endpoints and the intermediate value theorem give the complete maximizer characterization. Divide the first-order equation by $D$ to obtain $p_0ke^{-ka}=\eta+(C/D)a$. Raising the right side pointwise moves its unique intersection with the decreasing left side to a smaller $a$; raising $p_0$ moves the intersection upward. Strictness holds while the root remains interior.
\end{proof}

A reduction in the bank component $p_0e^{-ka}$ occurs for every positive tightening, yet the optimal policy can be zero when leakage is sufficiently strong. Even when the full crisis probability falls, tightening beyond the optimum lowers welfare because the marginal crisis benefit no longer covers access costs.

\begin{proposition}[Group incidence and compensation]
Group $i$ gains if and only if
\[
 d_i\{p_0(1-e^{-ka})-\eta a\}\ge c_i a^2/2.
\]
If group utilities are quasilinear and a date-zero numeraire transfer changes lifetime utility dollar for dollar, the compensating transfer required to offset that group's policy loss is
\[
 T_i(a)=\frac{c_i a^2/2-d_i\{p_0(1-e^{-ka})-\eta a\}}{1-\beta}.
\]
This expression is signed: negative values are amounts the group could pay while retaining baseline utility. A positive aggregate weighted welfare gain does not imply $T_i(a)\le0$ for every group.
\end{proposition}
\begin{proof}
The group-specific per-period difference was computed in the preceding proof. Discounting gives its lifetime difference, and a compensating date-zero transfer is its negative. To see the final assertion, take two groups: one has $d_1>0,c_1=0$, and the other has $d_2=0,c_2>0$. At any policy that lowers crisis probability, the first gains and the second loses. Choose its cost sufficiently small or its welfare weight sufficiently small for the weighted sum to be positive.
\end{proof}

The compensation formula does not assert that compensation is feasible. Transfers require a fiscal budget and may change behavior or financing costs. In particular, the sum of unweighted transfers is a different object from the weighted welfare objective.

\section*{Sharp ambiguity about the future}
Suppose a causal experiment or credible design identifies the expected discounted utility difference through date $H$,
\[
 D_H=\sum_{t=0}^H\beta^t\E[U_t(1)-U_t(0)],
 \qquad U_t(a)=\sum_i\omega_i u_i(c_{it}(a)).
\]
Assume only that $U_t(a)\in[u_-,u_+]$ almost surely for both regimes and all future dates, with $u_-<u_+$, and allow arbitrary future laws consistent with these bounds and the observed prefix.

\begin{theorem}[Sharp long-run welfare interval]
The identified set for the infinite-horizon welfare difference is exactly
\[
 \left[D_H-\frac{\beta^{H+1}(u_+-u_-)}{1-\beta},\quad
 D_H+\frac{\beta^{H+1}(u_+-u_-)}{1-\beta}\right].
\]
In particular, a positive measured prefix identifies a positive total effect only when it exceeds the displayed tail radius. This is a sharp statement for the declared unrestricted-continuation class, not for every structural macroeconomic model.
\end{theorem}
\begin{proof}
At each future date the utility difference lies in $[-(u_+-u_-),u_+-u_-]$. Summing the discounted bounds proves inclusion and absolute convergence. To attain the upper endpoint, keep the observed prefix law unchanged and set $U_t(1)=u_+$, $U_t(0)=u_-$ deterministically for all $t>H$. Reverse them for the lower endpoint. Both continuations satisfy the declared class. Mixtures of these two continuation systems, drawn after or independently of the fixed prefix, attain every intermediate value. Therefore the entire interval is sharp.
\end{proof}

If feasibility or equilibrium restrictions prohibit these continuations, the class must be narrowed before sharper claims are made. A finite-horizon risk estimate alone supplies no such restriction. Similarly, if a future crisis is defined by a fixed output-loss threshold, its probability and the utility severity $d_i$ are different primitives; one does not identify the other.

\section*{Normative sensitivity can change the optimum}
\begin{proposition}[A two-group weight threshold]
At a fixed policy with crisis-probability reduction $b>0$, let group 1 have gain $g_1>0$ and group 2 have loss $g_2<0$ per period, calculated from their access costs and crisis severities above. With weight $\omega$ on group 1, the policy raises discounted welfare exactly when
\[
 \omega>\frac{-g_2}{g_1-g_2},
\]
with equality giving zero gain. Both group effects and the crisis probability can be held fixed while the social ranking reverses.
\end{proposition}
\begin{proof}
The discounted gain is $[\omega g_1+(1-\omega)g_2]/(1-\beta)$. The denominator is positive and the numerator has positive slope $g_1-g_2$, so solving its strict positivity gives the threshold. No observable causal law changes when only the normative weight changes.
\end{proof}

\section*{Remaining identification and model gap}
A full evaluation must identify the total probability response including nonbank and foreign leakage, group consumption effects, and the long-run continuation. Endogenous policy tightening during latent risk buildups invalidates an unadjusted before--after interpretation. The exact optimum above presumes a stable response curve and does not solve this selection problem. It also excludes household insurance, endogenous investment and policy interactions. These are economically material limits rather than small technical extensions.

The IMF primary eLibrary text and publication metadata were checked. Chapters 3--4 distinguish real-activity benefits and costs, nonlinear effects and leakage; Chapter 6 discusses remaining research. The seven-author citation follows the report rather than the abbreviated landing-page author list. The constructed response curve and its welfare optimum are not empirical findings of that report.

\section{Completion Estimate}
30\%

This is a post hoc, heuristic estimate for the full catalogue target.
The attempt solves a heterogeneous-household welfare benchmark with leakage and derives sharp long-run utility-tail bounds plus group compensation and welfare-weight thresholds. It does not identify the policy response or long-run consumption effects under regulator selection, richer equilibrium leakage, insurance, investment, and interacting rules.

\begin{thebibliography}{9}
\bibitem{imf} N. Biljanovska, S. Chen, G. Gelos, D. Igan, M. S. Mart\'inez Per\'ia, E. Nier and F. Valencia (2023), \emph{Macroprudential Policy Effects: Evidence and Open Questions}, IMF Departmental Paper 2023/002. \url{https://www.elibrary.imf.org/view/journals/087/2023/002/article-A001-en.xml}. DOI: \url{https://doi.org/10.5089/9798400226304.087}.
\end{thebibliography}

\end{document}
